It is enough to verify the first equality.  Since the two sides are
smooth \(S\)-subgroups of \(G\), it is enough to check that they have
the same geometric points.  In the fibers over points \(s\ne s_0\),
this follows because \(\mathrm{PGL}_{2,\kappa(s)}\) is reductive and
\(T_s\) is a maximal torus in it for dimensional reasons
\textup{(cf.\ \SGARefLink{sga3:xix:lemma:1-11}{1.11})}.  In the fiber over \(s_0\), the verification is
immediate.  In particular, \(T\) is a maximal torus of both \(G\)
and \(G^0\).

\par\medskip\noindent
\textbf{5.12.}---
The section \(w\) of \(G\) defined in \SGARefLink{sga3:xix:lemma:5-10}{Lemma 5.10} normalizes \(T\).
It follows immediately \textup{(cf.\ \SGARefLink{sga3:xix:remark:2-4}{2.4})} that the Weyl group of
\(G\) is isomorphic to \((\mathbb Z/2\mathbb Z)_S\), and is therefore
finite over \(S\):
\[
  W_G(T)=\operatorname{Norm}_G(T)/T
        =(\mathbb Z/2\mathbb Z)_S.
\]
By contrast, \(W_{G^0}(T)\) is not finite over \(S\): it ``lacks one
point'' above \(s_0\).
\SGARefTarget{sga3:xix:paragraph:5:12}

\par\medskip\noindent
\textbf{\SGAInlineRefTarget{sga3:xix:paragraph:5-13}5.13.}---
The open immersion \(G^0\to G\) is not a closed immersion, since
\(G^0\) is dense in \(G\); it is nevertheless an affine morphism
\textup{(and hence \(G^0\) is affine over \(S\))}.  Indeed,
\(G^0_{s_0}\) is closed in \(G_{s_0}\), which is closed in \(G\), so
its complement \(U\) in \(G\) is open.  The open subsets \(G^0\) and
\(U\) cover \(G\), and it is enough to verify that the immersions
\[
  G^0\longrightarrow G^0
  \qquad\text{and}\qquad
  G^0\cap U\longrightarrow U
\]
are affine.  This is trivial for the first; for the second, observe
that \(U\cap G^0\) is defined in \(U\) by the condition \(b\ne0\).

We have thus constructed a smooth affine \(S\)-group \(G^0\), with
connected fibers, possessing a self-centralizing maximal torus \(T\)
whose Weyl group \(W_{G^0}(T)\) is not finite \textup{(compare
\SGARefLink{sga3:xix:theorem:2:5}{Theorem 2.5})}.
\SGARefTarget{sga3:xix:paragraph:5:13}

\SGANextSectionTarget{sga3:xix:section:6}\SGASection{6}{Local existence of maximal tori.  The Weyl group}

In the proof of \SGARefLink{sga3:xix:theorem:2:5}{Theorem 2.5} we used a result from Exposé XI on the
local existence, for the étale topology, of maximal tori.  The proof
in Exposé XI uses a deep representability result \textup{(\SGARefLink{sga3:xi:theorem:4:1}{XI 4.1})}.
In the particular case under consideration, one can give another,
far more elementary proof, based on the ideas of Exposé XII, no.~7.

\par\medskip\noindent
\SGAInlineRefTarget{sga3:xix:proposition:6-1}\textbf{\textit{\SGARefLink{sga3:xix:proposition:6:1}{Proposition 6.1}.}}---
\textit{Let \(S\) be a scheme, let \(G\) be a smooth affine
\(S\)-group scheme with connected fibers over \(S\), and let
\(s_0\in S\) be a point such that the maximal tori of the geometric
fiber \(G_{\bar s_0}\) are self-centralizing.  There exist an étale
morphism \(S'\to S\) whose image contains \(s_0\) and a split maximal
torus \(T\) of \(G_{S'}\).}
\SGARefTarget{sga3:xix:proposition:6:1}

\SGAProof{Proof}
First, we may suppose that \(S\) is affine.
\footnote{Editor's note: the preceding sentence and the reference
to EGA IV below have been added.}
Since \(G\) is of finite presentation over \(S\), we may successively
suppose that \(S\) is noetherian, local, and then henselian with
separably closed residue field \textup{(cf.\ EGA IV, 8.12,
\S\S 8.8 and 18.8)}.  Thus let
\[
  S=\operatorname{Spec}(A),
  \qquad
  k=\kappa(s_0),
\]
where \(A\) is henselian and \(k\) is separably closed.  Choose a
maximal torus \(T_0\) of \(G_0=G_k\).  Such a torus exists, for
example because the scheme of maximal tori of \(G_0\) is smooth over
\(k\) \textup{(\SGARefLink{sga3:xii:7.1:item:c}{Exposé XII, 7.1(c)})}.  Since \(k\) is separably
closed, \(T_0\) is split \textup{(cf.\ \SGARefLink{sga3:x:proposition:1:4}{Exposé X, 1.4})}, and is
therefore given by a group monomorphism
\[
  f_0\colon \mathbb G_{m,k}^{\,r}\longrightarrow G_0.
\]
